\documentclass[aspectratio=169]{beamer}

\usetheme{metropolis}

\usepackage{graphicx}
\usepackage{booktabs}
\usepackage{hyperref}

\title{Optimization-Based Decision Support for Equitable FQHC Expansion in Chicago}
\subtitle{Meeting with Prof. Sage Kim}
\author{Alaittin Kirtisoglu \and Hemanshu Kaul}
\institute{Department of Applied Mathematics \\ Illinois Institute of Technology}
\date{\today}

\begin{document}

\maketitle

% ---------------------------------------------------------------
\begin{frame}{The Problem: Four Gaps in the HRSA System}

\begin{enumerate}
  \item<1-> \textbf{No location evaluation tools.} When an organization applies
        for a New Access Point grant, there is no tool to rigorously answer
        \emph{why this location}.

  \item<2-> \textbf{The IMU score is outdated.} The formula dates from the
        1970s, omits insurance coverage, transportation, and chronic disease
        burden, and uses a hard binary cutoff.

  \item<3-> \textbf{The MUA designation process is uneven.} 87 census tracts in
        Chicago meet IMU eligibility but were never designated, predominantly in
        Black South Side neighborhoods (Kim et al., 2020).

  \item<4-> \textbf{No referral network analysis.} HRSA mandates referral
        relationships with hospitals and specialists, but no one has mapped how
        referrals actually flow in Chicago.
\end{enumerate}

\end{frame}

% ---------------------------------------------------------------
\begin{frame}{The Optimization Model in Plain English}

\textbf{Inputs:} facility locations, demand, capacity, travel constraints.

\medskip
\textbf{Output:} many high-quality configurations for decision-makers to
\emph{compare and choose from} --- not a single optimal solution.

\bigskip
\textbf{Two hierarchical levels:}

\begin{itemize}
  \item \textbf{Districts} --- catchment areas around each primary care
        facility. Addresses HRSA's catchment area design requirement.
  \item \textbf{Superdistricts} --- groups of primary care districts sharing a
        referral hospital. Addresses HRSA's referral relationship requirement.
\end{itemize}

\bigskip
The framework is built to respect equity, travel time by public transit, and
user-defined need measures.

\end{frame}

% ---------------------------------------------------------------
\begin{frame}{Near-Term Paper: FQHC Optimization with UDS Projection}

\textbf{Data gap.} UDS data (visits, FTEs) is reported at the
\emph{organization} level, not per site. 25 parent organizations operate 166
FQHC sites in Chicago.

\bigskip
\textbf{Our approach.} Project aggregate organization-level UDS onto individual
sites proportionally --- by provider count, facility size, or reported FTEs.

\bigskip
\textbf{What this enables.}
\begin{itemize}
  \item Run the full FQHC optimization framework with realistic demand and
        capacity estimates.
  \item Publish the methodology now, without waiting on a data use agreement.
  \item Demonstrate the analytical value of having accurate site-level data.
\end{itemize}

\bigskip
\textbf{Honest caveat.} This is an approximation. Its purpose is to motivate
the case for true site-level data.

\end{frame}

% ---------------------------------------------------------------
\begin{frame}{Follow-up: Chicago Physician Referral Network}

\textbf{Data source.} CMS Physician Shared Patient Patterns --- physicians are
nodes, shared Medicare patients within 30-day windows form edges.

\bigskip
\textbf{What we would build.}
\begin{itemize}
  \item A Chicago-specific physician referral network.
  \item Overlay with FQHC and hospital data to identify \emph{referral deserts}
        --- communities where patients have no viable path from primary care to
        specialist care.
\end{itemize}

\bigskip
\textbf{Connection to your CBO capacity work.} Referral pathways combined with
real-time capacity data is what turns a static map into a dynamic routing tool.

\bigskip
\textbf{Data access.} ResDAC application, or a data use agreement through UIC.

\end{frame}

% ---------------------------------------------------------------
\begin{frame}{What We Offer and What We Need}

\textbf{What we bring:}
\begin{itemize}
  \item A spatial optimization framework for hierarchical facility location and
        districting.
  \item Mathematical methods for equitable resource allocation.
  \item An interactive decision-support dashboard integrating 166 FQHC sites,
        128 hospitals, 61{,}460 providers, and 33 health indicators.
\end{itemize}

\bigskip
\textbf{What we need:}
\begin{itemize}
  \item Health domain expertise to ground the modeling choices.
  \item Data access through UIC (site-level UDS, CMS referral data).
  \item Understanding of what decision-makers actually need from such tools.
\end{itemize}

\end{frame}

\end{document}
